% Part: incompleteness
% Chapter: representability-in-q
% Section: representable-comp

\documentclass[../../../include/open-logic-section]{subfiles}

\begin{document}

\olfileid[fr]{inc}{req}{rpc}
\olsection{Les fonctions représentables dans $\Th{Q}$ sont calculables}

Nous allons démontrer que toute fonction représentable dans~$\Th{Q}$
est calculable. Il nous faut d'abord établir un lemme sur ces
fonctions.

\begin{lem}\ollabel{lem:rep-q}
  Si $f(x_0,\dots,x_k)$ est représentable dans~$\Th{Q}$, il existe
  une !!{formula}~$!A_f(x_0,\dots,x_k,y)$ telle que
  \[
  \Th{Q} \Proves !A_f(\num{n_0}, \dots, \num{n_k}, \num{m})
  \quad\text{ssi}\quad m = f(n_0, \dots, n_k).
  \]
\end{lem}

\begin{proof}
  L'implication « si » découle de la
  \olref[int]{defn:representable-fn}\olref[int]{defn:rep:a}.
  Pour l'implication réciproque, supposons que
  $\Th{Q} \Proves !A_f(\num{n_0},\dots,\num{n_k},\num{m})$
  mais que $m\neq f(n_0,\dots,n_k)$, et posons
  $l=f(n_0,\dots,n_k)$. Par la
  \olref[int]{defn:representable-fn}\olref[int]{defn:rep:a},
  $\Th{Q} \Proves !A_f(\num{n_0},\dots,\num{n_k},\num{l})$.
  D'après la \olref[int]{defn:representable-fn}\olref[int]{defn:rep:b},
  $\Th{Q}$ démontre aussi
  $\lforall[y][(!A_f(\num{n_0},\dots,\num{n_k},y)\lif\num{l}=y)]$.
  En appliquant les règles logiques à cet énoncé et à l'hypothèse
  $\Th{Q} \Proves !A_f(\num{n_0},\dots,\num{n_k},\num{m})$,
  on obtient $\Th{Q} \Proves \eq[\num{l}][\num{m}]$.
  D'autre part, \olref[bre]{lem:q-proves-neq} donne
  $\Th{Q} \Proves \eq/[\num{l}][\num{m}]$. La théorie $\Th{Q}$
  serait donc incohérente. C'est impossible : le modèle standard
  satisfait $\Th{Q}$ (voir la \olref[int][def]{def:standard-model}),
  soit $\Sat{N}{\Th{Q}}$, et toute théorie satisfaisable est
  cohérente par le théorème de correction
  (\tagrefs{prfAX/{fol:axd:sou:cor:consistency-soundness},
prfSC/{fol:seq:sou:cor:consistency-soundness},
prfND/{fol:ntd:sou:cor:consistency-soundness},
prfTab/{fol:tab:sou:cor:consistency-soundness}}).
\end{proof}

\begin{lem}
Toute fonction représentable dans $\Th{Q}$ est calculable.
\end{lem}

\begin{proof}
Commençons par l'idée intuitive. Pour calculer~$f$, énumérons
toutes les !!{derivation}s possibles dans le langage de
l'arithmétique; cette énumération peut se faire mécaniquement.
Pour chacune, vérifions si elle démontre, à partir des axiomes
de~$\Th{Q}$, !!a{formula} de la forme
$!A_f(\num{n_0},\dots,\num{n_k},\num{m})$, où $!A_f$ représente
$f$ comme dans le \olref{lem:rep-q}. Si oui, alors
$m=f(n_0,\dots,n_k)$ d'après le \olref{lem:rep-q}, et nous avons
trouvé la valeur de~$f$. La recherche s'arrête, car
$\Th{Q} \Proves !A_f(\num{n_0},\dots,\num{n_k},
\num{f(n_0,\dots,n_k)})$; nous rencontrerons donc une
!!{derivation} appropriée.

Cette description n'est pas encore tout à fait précise : notre
procédure manipule des !!{derivation}s et des !!{formula}s, et
non seulement des nombres, et nous n'avons pas expliqué pourquoi
on peut énumérer mécaniquement toutes les !!{derivation}s
possibles. Mais nous avons vu comment coder les termes, les
!!{formula}s et les !!{derivation}s par des numéros de G\"odel.
Nous avons aussi défini un modèle précis de calcul, celui des
fonctions récursives générales. La relation
$\Prf[\Th{Q}](d,y)$, qui vérifie qu'un code~$d$ est un témoin de
preuve, à partir des axiomes de~$\Th{Q}$, de la !!{formula} de
numéro de G\"odel~$y$, est récursive primitive. Ici, les
prémisses de ce test sont les huit axiomes de~$\Th{Q}$, et non
l'ensemble déjà clos de tous ses théorèmes. Les autres fonctions
récursives primitives dont nous aurons besoin sont $\fn{num}$
(\olref[art][trm]{prop:num-primrec}) et $\fn{Subst}$
(\olref[art][sub]{prop:subst-primrec}). Elles permettront de
définir~$f$ par minimisation; $f$ est donc récursive.

Définissons d'abord
\begin{multline*}
  A(n_0, \dots, n_k, m) = \\
  \fn{Subst}(\fn{Subst}(\dots\fn{Subst}(\Gn{!A_f}, \fn{num}(n_0), \Gn{x_0}),\\ \dots),
  \fn{num}(n_k),  \Gn{x_k}), \fn{num}(m), \Gn{y})
\end{multline*}
L'expression est compliquée, mais cette fonction donne simplement
$A(n_0,\dots,n_k,m)=\Gn{!A_f(\num{n_0},\dots,\num{n_k},\num{m})}$.

Considérons maintenant la relation $R(n_0,\dots,n_k,s)$, vraie si
$s$ code un couple dont la première composante $(s)_0$ est un
témoin de preuve, à partir des axiomes de~$\Th{Q}$, de
$!A_f(\num{n_0},\dots,\num{n_k},\num{(s)_1})$. On pose
\begin{multline*}
R(n_0,\dots,n_k,s) \quad\text{ssi}\quad
\fn{SeqCode}(s)\land \len{s}=2\\
\land \Prf[\Th{Q}]((s)_0,A(n_0,\dots,n_k,(s)_1)).
\end{multline*}
Si nous trouvons un~$s$ vérifiant $R(n_0,\dots,n_k,s)$, nous
obtenons la paire $((s)_0,(s)_1)$ : son premier élément code un
témoin de preuve de
$!A_f(\num{n_0},\dots,\num{n_k},\num{(s)_1})$ à partir
des axiomes de~$\Th{Q}$, et son second élément est la valeur
cherchée. Chercher~$s$ revient donc à chercher les nombres $d$
et $m$ de l'argument informel. Une fonction calculable qui trouve
un tel~$s$ se définit par minimisation régulière. En effet,
$R$ possède toujours un témoin : pour tous $n_0,\dots,n_k$,
il existe un témoin de preuve~$\delta$ attestant
$\Th{Q} \Proves !A_f(\num{n_0},\dots,\num{n_k},
\num{f(n_0,\dots,n_k)})$, de sorte que $R(n_0,\dots,n_k,s)$
est vraie pour
$s=\tuple{\Gn{\delta},f(n_0,\dots,n_k)}$.
Nous pouvons donc écrire
\[
f(n_0,\dots,n_{k}) = (\umin{s}{R(n_0, \dots, n_k, s)})_1.
\]
\end{proof}

\begin{explain}
Le test d'un témoin utilise la liste finie des axiomes de
$\Th{Q}$ et les règles du calcul. C'est cette vérification
de preuves qui permet de chercher la valeur de~$f$ sans disposer
d'une procédure de décision pour les théorèmes de~$\Th{Q}$.
\end{explain}

\end{document}
